\documentclass[12pt]{article}
\usepackage[utf8]{inputenc}
\usepackage[slovene]{babel}
\usepackage{amsmath}
\usepackage{graphicx}
\usepackage{float}

\title{1. domača naloga pri predmetu Numerična matematika}
\author{Anja Abramovič}
\date{}

\begin{document}
\maketitle

\section{Opis naloge}
Naloga obravnava reševanje sistema linearnih enačb $A\mathbf{x} = \mathbf{b}$, kjer je $A$ diagonalno dominantna redka matrika. Implementirali smo ustrezno podatkovno strukturo za redke matrike ter iterativno metodo SOR za reševanje sistema. Metodo smo uporabili za vložitev grafov v ravnino s fizikalno metodo.

\section{Redka matrika}
Redka matrika je matrika, pri kateri je velika večina elementov ničelnih. Shranjevanje celotne $n \times n$ matrike bi zahtevalo $\mathcal{O}(n^2)$ prostora, čeprav večina elementov ne nosi nobene informacije. Zato definiramo podatkovno strukturo \texttt{RedkaMatrika}, ki hrani le neničelne elemente. Za vsako vrstico $i$ sta shranjeni dve tabeli:
\begin{itemize}
    \item $V[i]$ — neničelne vrednosti v vrstici $i$,
    \item $I[i]$ — indeksi stolpcev teh vrednosti.
\end{itemize}
Torej velja $V[i][j] = a_{i, I[i][j]}$. Prostorska zahtevnost se tako zmanjša na 
$\mathcal{O}(k)$, kjer je $k$ število neničelnih elementov.

Na primer, matriko
\begin{equation}
    A = \begin{pmatrix} 4 & 0 & 2 \\ 0 & 3 & 0 \\ 4 & 0 & 5 \end{pmatrix}
\end{equation}
shranimo kot:
\begin{align*}
    V &= \big([4,\ 2],\ [3],\ [4,\ 5]\big), \\
    I &= \big([1,\ 3],\ [2],\ [1,\ 3]\big).
\end{align*}

\noindent Za podatkovno strukturo so implementirane naslednje operacije:
\begin{itemize}
    \item \texttt{getindex}: vrne element $a_{ij}$, oziroma $0$, če element ni shranjen,
    \item \texttt{setindex!}: nastavi element $a_{ij}$ oziroma doda nov neničelni element,
    \item \texttt{size}: vrne velikost matrike,
    \item \texttt{*}: pomnoži matriko z vektorjem $A\mathbf{x}$, pri čemer upošteva le neničelne elemente,
    \item \texttt{-}: vrne negacijo matrike $-A$.
\end{itemize}

\section{Metoda SOR}
Za reševanje sistema $A\mathbf{x} = \mathbf{b}$ uporabimo iterativno metodo SOR 
(Successive Over-Relaxation). V vsakem koraku posodobimo posamezne komponente 
rešitve po formuli:
\begin{equation}
    x_i^{(k+1)} = (1 - \omega) x_i^{(k)} + \frac{\omega}{a_{ii}} 
    \left( b_i - \sum_{j < i} a_{ij} x_j^{(k+1)} - \sum_{j > i} a_{ij} x_j^{(k)} \right),
\end{equation}
kjer je $\omega \in (0, 2)$ parameter sprostitve. Iteracija se ustavi, ko je norma ostanka dovolj majhna:
\begin{equation}
    \|A\mathbf{x}^{(k)} - \mathbf{b}\|_\infty < \delta.
\end{equation}
Hitrost konvergence je močno odvisna od izbire $\omega$, optimalna vrednost 
minimizira število potrebnih iteracij.

Pri uporabi redke matrike posodobitev $x_i^{(k+1)}$ izračunamo le iz neničelnih 
elementov vrstice $i$, kar zmanjša časovno zahtevnost posamezne iteracije z 
$\mathcal{O}(n)$ na $\mathcal{O}(k)$, kjer je $k$ povprečno število neničelnih 
elementov v vrstici.

\section{Rezultati}
Metodo SOR smo uporabili za vložitev grafov v ravnino s fizikalno metodo, ki temelji na 
modelu harmoničnih vzmeti. Koordinate vozlišč določimo kot ravnovesje sil, kar vodi do 
sistema linearnih enačb z redko, diagonalno dominantno in negativno definitno matriko.

Metodo smo preizkusili na grafu krožne lestve z 16 vozlišči. Zunanja vozlišča so fiksirana
enakomerno na enotski krožnici, notranja pa postavimo v ravnovesje s pomočjo metode SOR 
(slika~\ref{fig:krozna}). Algoritem preizkusimo še na pravokotni mreži $6 \times 6$, 
kjer fiksiramo le štiri vogalna vozlišča (slika~\ref{fig:mreza}).

\begin{figure}[H]
    \centering
    \includegraphics[width=0.6\textwidth]{krozna.png}
    \caption{Krožna lestev z 16 vozlišči, vložena s fizikalno metodo.}
    \label{fig:krozna}
\end{figure}

\begin{figure}[H]
    \centering
    \includegraphics[width=0.6\textwidth]{mreza.png}
    \caption{Pravokotna mreža $6 \times 6$, vložena s fizikalno metodo. Pritrjeni so le vogali.}
    \label{fig:mreza}
\end{figure}

Hitrost konvergence metode SOR je močno odvisna od izbire parametra $\omega$. Na 
sliki~\ref{fig:konvergenca} je prikazano število iteracij v odvisnosti od $\omega$ 
za sistem krožne lestve. Optimalna vrednost $\omega^* \approx 1.12$ minimizira 
število potrebnih iteracij na 21 iteracij.

\begin{figure}[H]
    \centering
    \includegraphics[width=0.6\textwidth]{konvergenca.PNG}
    \caption{Odvisnost števila iteracij metode SOR od parametra $\omega$ za sistem krožne lestve.}
    \label{fig:konvergenca}
\end{figure}

\end{document}